% ==========================================================================
% PlanPilot 项目经历（改写版）
%
% 关键点已加粗。
% \ResumeItem 的可选参数必须是 [..] 且紧贴命令，\ResumeItem{标题}[副标题][右对齐]
% —— 中间不能有换行，否则 [ ] 会被原样打印成方括号。
% ==========================================================================

\ResumeItem{PlanPilot — 本地优先的 LLM Agent 日程规划系统}[FastAPI / SQLAlchemy / OR-Tools CP-SAT / Function Calling]
\begin{itemize}
\item \textbf{问题背景}: 直接由 LLM 输出日程存在时间计算幻觉，且无法保证冲突检测与排期的严谨性，难以支撑真实场景的多约束排期。
\item \textbf{架构解耦}: 设计"\textbf{LLM 意图理解 + 确定性规划引擎}"分离架构。LLM 仅负责对话理解、工具选择与多步推理；时间计算交由本地纯 Python 引擎与 \textbf{CP-SAT 求解器}执行，引擎不访问数据库、不调用模型，\textbf{同输入必同输出}，模型不可用时排期能力不受影响。
\item \textbf{Tool Calling 框架}: 从零搭建 Tool Registry 与 Agent Loop，封装 \textbf{37 个工具}、\textbf{10 个意图路由 Skill}；实现上下文动态裁剪、多轮对话记忆、结构化 Observation 回传与工具级白名单隔离。
\item \textbf{可靠性工程}: 针对 Agent 死循环与模型侧故障，设计步数/调用预算、重复调用拦截、危险操作二次确认三层阻断；将 LLM 异常细分为鉴权失败/限流/瞬时故障三类并施加\textbf{指数退避}（60s→600s 带抖动），配合环境代理失效直连回退与\textbf{补偿式事务}（多步执行失败整体回滚，\textbf{6 张表}快照还原）。
\item \textbf{业务约束求解}: 集成高德地图路线估算校验相邻行程真实可达性，识别"\textbf{时间不重叠但路上赶不及}"的隐性冲突；支持多段行程组原子创建/顺延/删除、计划版本结构化 diff 与一键撤销、只读式风险分析与\textbf{多方案对比决策}。
\item \textbf{工程化与可观测}: 基于 FastAPI + SQLAlchemy 构建 RESTful API，统一 UTC 存储与本地时区渲染；实现持久化执行轨迹（分阶段耗时、刷新可回看）与隔离数据库的 Agent 端到端评测（\textbf{5 个场景}，统计工具选择正确率/约束通过率/确认安全性）；编写 \textbf{220 个}自动化测试用例，核心逻辑回归通过率 \textbf{100\%}。
\end{itemize}

% ==========================================================================
% 下面两段是原简历中 \ResumeItem 参数形式修正
% 原写法：\ResumeItem{标题}\n[副标题]\n[日期]  ->  方括号被原样渲染
% 修正：把可选参数紧贴命令，同一行内给出
% ==========================================================================

% ---------- 教育背景 ----------
\ResumeItem{南京航空航天大学}[\textnormal{计算机科学与技术 | 硕士}][2025.09—2028.03]
\ResumeItem{南京航空航天大学}[\textnormal{人工智能（培优班） | 本科}][2021.09—2025.06]

% ---------- 科研经历 ----------
\ResumeItem{BioBridge — 基于生物视觉先验的 EEG-图像跨模态解码}[PyTorch / 跨模态对比学习]
\begin{itemize}
\item \textbf{问题背景}: 脑电（EEG）信号与图像表征存在巨大模态差异，高效精准的零样本跨模态对齐是脑机接口领域的核心难题。
\item \textbf{算法设计}: 使用 PyTorch 设计视觉先验引导的分层特征提取网络，引入多级双向对比学习模块，缩小 EEG 信号与视觉特征的语义鸿沟。
\item \textbf{实验成果}: 在 THINGS-EEG 200 类零样本图像检索任务中取得 Top-1 \textbf{80.5\%}、Top-5 \textbf{97.6\%} 的准确率；成果在投 ICLR 2027。
\end{itemize}
